\documentclass[12pt]{article}
\usepackage{amsmath}
\usepackage{geometry} % Required for customizing margins
\geometry{margin=1in} % Sets 1-inch margins on all sides
\begin{document}

\title{CSCE 350 Section 001: Assignment \#3}
\author{Jacob C. Vaught}
\maketitle

\section*{Problem \#1[42 Points]}
Solve the following recurrence relations using the method of backward substitutions. Give the solution to the problem in terms of `n'.

\begin{enumerate}
    \item \( X(n) = X(n-1) + 5 \) for \( n > 0 \), \( X(0)=1 \)
    \item \( X(n) = X(n-1) + 5n \) for \( n > 0 \), \( X(1)=5 \)
    \item \( X(n) = X\left(\frac{n}{4}\right) + n \) for \( n > 1 \), \( X(1)=1 \)
\end{enumerate}

\section*{Solutions}

\subsection*{Part a[12 Points]}
For the recurrence relation \( X(n) = X(n-1) + 5 \) with \( X(0) = 1 \):

\subsubsection*{Backward Substitutions}
Starting with \(X(n)\):
\begin{align*}
X(n) &= X(n-1) + 5 \\
X(n-1) &= X(n-2) + 5 \\
X(n-2) &= X(n-3) + 5 \\
&\vdots \\
X(2) &= X(1) + 5 \\
X(1) &= X(0) + 5
\end{align*}

\subsubsection*{Final Solution}
Substituting backward and collapsing:
\begin{align*}
X(n) &= (X(n-1) + 5) \\
&= [(X(n-2) + 5) + 5] \\
&= [(X(n-3) + 5) + 5 + 5] \\
&\vdots \\
&= 5n + X(0) \\
&= 5n + 1
\end{align*}

\subsection*{Part b[15 Points]}
For the recurrence relation \( X(n) = X(n-1) + 5n \) with \( X(1) = 5 \):

\subsubsection*{Backward Substitutions}
Starting with \(X(n)\):
\begin{align*}
X(n) &= X(n-1) + 5n \\
X(n-1) &= X(n-2) + 5(n-1) \\
X(n-2) &= X(n-3) + 5(n-2) \\
&\vdots \\
X(2) &= X(1) + 5 \cdot 2 \\
X(1) &= 5
\end{align*}

\subsubsection*{Final Solution}
Substituting backward and collapsing:
\begin{align*}
X(n) &= (X(n-1) + 5n) \\
&= [(X(n-2) + 5(n-1)) + 5n] \\
&= [(X(n-3) + 5(n-2)) + 5(n-1) + 5n] \\
&\vdots \\
&= 5 + 5 \cdot 2 + 5 \cdot 3 + \cdots + 5n \\
&= 5 + 5 \cdot \sum_{i=1}^{n} i \\
&= 5 + 5 \cdot \frac{n(n + 1)}{2} \\
&= 5 + \frac{5n(n + 1)}{2}
\end{align*}

\subsection*{Part c[15 Points]}
For the recurrence relation \( X(n) = X\left(\frac{n}{4}\right) + n \) with \( X(1) = 1 \), determine \( X(n) \) using the backwards substitution method.

\subsubsection*{Backward Substitution}
We'll start by writing down the recurrence relation:

\begin{align*}
X(n) &= X\left(\frac{n}{4}\right) + n
\end{align*}
Now, let's expand this recurrence relation using backwards substitution:

\begin{align*}
X(n) &= n + X\left(\frac{n}{4}\right) \\
&= n + \frac{n}{4} + X\left(\frac{n}{16}\right) \\
&= n + \frac{n}{4} + \frac{n}{16} + X\left(\frac{n}{64}\right) \\
&\vdots \\
&= n + \frac{n}{4} + \frac{n}{16} + \frac{n}{64} + \cdots + \frac{n}{4^k} + X\left(\frac{n}{4^{k+1}}\right)
\end{align*}
Since \( X(1) = 1 \), we will have \( X\left(\frac{n}{4^{k+1}}\right) = 1 \) when \( \frac{n}{4^{k+1}} = 1 \), or equivalently, \( k+1 = \log_4{n} \).

\subsubsection*{Closed-Form Expression}
Putting all these terms together, we have:

\[
X(n) = n \left(1 + \frac{1}{4} + \frac{1}{16} + \frac{1}{64} + \cdots + \frac{1}{4^k}\right) + 1
\]
\newline\newline
The geometric series \(1 + \frac{1}{4} + \frac{1}{16} + \cdots + \frac{1}{4^k}\) has \( a = 1 \) and \( r = \frac{1}{4} \). For \( k = \log_4{n} - 1 \), the sum \( S \) of this series can be simplified to:

\[
S = \frac{(4/4) - (1/4^{\log_4{n}})}{(4 - 1)/4}
\]
Therefore, we find:

\[
X(n) = n \times \frac{(4/4) - (1/4^{\log_4{n}})}{(4 - 1)/4} + 1
\]
Simplifying, we get:

\[
X(n) = \frac{4n - 1}{3}
\]


\newpage
\section*{Problem \#2 [20 Points]}
Solve the linear second-order recurrence relation \(X(n) - 3X(n-1) + X(n-2) = 3\) for \( n > 1 \) with \( X(0) = 0 \) and \( X(1) = 1 \). Give both the general and particular solutions.
\subsubsection*{Homogeneous Solution}
First, let's tackle the homogeneous part, meaning we look at the equation without the constant term on the right-hand side:
\[
X(n) - 3X(n-1) + X(n-2) = 0
\]
We form a characteristic equation from this homogeneous part:
\[
r^2 - 3r + 1 = 0
\]
Using the quadratic formula, we find the roots to be:
\[
r_1 = \frac{3 + \sqrt{5}}{2}, \quad r_2 = \frac{3 - \sqrt{5}}{2}
\]
So the general form of the homogeneous solution \( X_h(n) \) is:
\[
X_h(n) = A(r_1)^n + B(r_2)^n
\]
Here, \( A \) and \( B \) are constants we need to determine.

\subsubsection*{Applying Initial Conditions to Homogeneous Solution}
Using the initial conditions \( X(0) = 0 \) and \( X(1) = 1 \), we can find \( A \) and \( B \):
\[
0 = A(r_1)^0 + B(r_2)^0 = A + B
\]
\[
1 = A(r_1)^1 + B(r_2)^1 = A \frac{3 + \sqrt{5}}{2} + B \frac{3 - \sqrt{5}}{2}
\]
Solving these equations yields \( A = 1 \) and \( B = -1 \).

\subsubsection*{Particular Solution}
The non-homogeneous term (the constant 3 on the right-hand side of the original equation) indicates that we should look for a constant particular solution \( X_p(n) \). Let it be \( C \):
\[
C - 3C + C = 3 \implies C = -\frac{3}{2}
\]

\subsubsection*{General Solution}
Finally, we combine the homogeneous and particular solutions to form the general solution:
\[
X(n) = X_h(n) + X_p(n)
\]
\[
X(n) = \left( \frac{3 + \sqrt{5}}{2} \right)^n - \left( \frac{3 - \sqrt{5}}{2} \right)^n - \frac{3}{2}
\]
This gives us a closed-form expression for \( X(n) \), answering the original question.

\newpage

\section*{Problem \#3[38 Points]}
Consider the algorithm MinDistance for finding the minimum distance between two elements in an array \(A[0..n-1]\).

\subsection*{Algorithm Description}
\begin{verbatim}
Algorithm MinDistance(A[0..n - 1])
    dmin ← \inf
    for i ← 0 to n - 1 do
        for j ← 0 to n - 1 do
            if i != j and sqrt((A[i]-A[j])^2) < dmin
                dmin ← sqrt((A[i]-A[j])^2)
    return dmin
\end{verbatim}

\subsection*{a) Basic Operation[3 Points]}
The basic operation in this algorithm is the computation of \( \sqrt{(A[i]-A[j])^2} \).

\subsection*{b) Worst and Best Case[8 Points]}
\subsubsection*{Worst Case}
The worst-case scenario occurs when every possible pair \((i, j)\) has to be checked, resulting in a total of \(n \times (n-1) = n^2 - n\) basic operations.
\subsubsection*{Best Case}
The best-case scenario is the same as the worst case in this algorithm, \(n^2 - n\), because all pairs must be checked regardless of the array's elements.

\subsection*{c) Algorithm Improvements[27 Points]}

\subsubsection*{i) Pseudo-Code of New Algorithm[10 Points]}
\begin{verbatim}
Algorithm ImprovedMinDistance(A[0..n - 1])
    Sort A in ascending order
    dmin ← \inf
    for i ← 0 to n - 2 do
        temp ← A[i+1] - A[i]
        if temp < dmin
            dmin ← temp
    return dmin
\end{verbatim}

\subsubsection*{ii) Description of Changes[7 Points]}
\begin{enumerate}
    \item The array \(A\) is sorted in ascending order initially. Sorting ensures that the closest elements will be adjacent, thereby eliminating the need for nested loops.
    \item A single for-loop is used to find the minimum distance between adjacent elements, which replaces the original nested for-loops.
    \item The square root and square operations from the original algorithm (\( \sqrt{(A[i]-A[j])^2} \)) have been removed. They are replaced by the simpler \( A[i+1] - A[i] \), which implicitly serves as the absolute difference in the sorted array.
\end{enumerate}

These changes improve the algorithm by reducing its time complexity. While the original algorithm had a time complexity of \(O(n^2)\), the new algorithm reduces it to \(O(n \log n)\) due to sorting.

\subsubsection*{iii) Basic Operation and Count in New Algorithm[10 Points]}
The new basic operation in the improved algorithm is \( \text{temp} \leftarrow A[i+1] - A[i] \). This operation directly contributes to finding the smallest distance between elements in the array.
\newline\newline
For an array of size \(n\), this basic operation is performed \(n-1\) times. When combined with the \(O(n \log n)\) time complexity for sorting, the total time complexity becomes \(O(n \log n + n) = O(n \log n)\).


\end{document}
